\documentclass[10pt,a4paper]{amsart}
\usepackage[margin=19mm]{geometry}
\usepackage{amsmath,amssymb,mathtools}
\usepackage[scr=rsfs,cal=euler]{mathalfa}
\usepackage{xfrac}
\usepackage[unicode,hidelinks,bookmarks=false]{hyperref}
\newcommand{\ie}{i.e.\ }
\newcommand{\cf}{cf.\ }
\newcommand{\resp}{respectively\ }
\newcommand{\Subsection}{\subsection}
\providecommand{\nobreakdash}{\nobreak\hbox{-}}
\setlength{\emergencystretch}{2em}
\multlinegap0pt
\usepackage{amssymb}
\usepackage{mathtools}
\usepackage[shortlabels]{enumitem}
\newtheorem{lemma}{Lemma}
\newtheorem{corollary}{Corollary}
\newtheorem{theorem}{Theorem}
\newcommand{\F}{\mathbb{F}}
\title{Sharp Asymptotics for Abelian Covers of Groups with Bounded Noncommutativity}
\author{Guillaume Lecomte}
\date{Source: arXiv:2608.20507v1 (CC BY 4.0)}
\begin{document}
\maketitle

\noindent\emph{Source and licence.} Sharp Asymptotics for Abelian Covers of Groups with Bounded Noncommutativity. Version arXiv:2608.20507v1 (CC BY 4.0), distributed by arXiv under the Creative Commons Attribution 4.0 International licence, http://creativecommons.org/licenses/by/4.0/, as recorded in the arXiv licence metadata for this exact version and shown on the abstract page https://arxiv.org/abs/2608.20507v1. Full licence notice: \texttt{licenses/arxiv-2608.20507-CC-BY-4.0.txt}.\par\smallskip

\noindent\textit{Editorial scope.} Counted unit: the article's theorem thm:pgroup (source0.tex lines 524--538). The article's complete original proof of it is delivered with it (source0.tex lines 540--581, one \texttt{proof} environment of 42 lines). No mathematics is written, repaired or re-derived: the statement and the proof are the article's own lines, in the article's own notation, and no retained line is substituted or re-flowed.  Retained uncounted as the source-local context the proof needs: lines 132--137; lines 143--148; lines 200--205; lines 301--303; lines 347--353; lines 448--453; lines 489--499.  Nothing was omitted from any retained span, so the package carries no omission record.  No retained line contains a citation, so the wrapper carries no bibliography and no bibliography entry.  No retained line contains a figure environment, an image-inclusion command or drawing code, so no image is delivered and the package declares no asset dependency.  Editorial formatting only, in addition: an AMS \texttt{amsart} preamble that restates the article's own declarations and macros used by the retained lines, namely the environments \texttt{lemma}, \texttt{corollary}, \texttt{theorem} on separate counters, together with the standard packages those lines need.  The retained environments print in this document as Lemma 1, Corollary 1, Theorem 1 in order of appearance, whereas the article prints the same items with the numbering of its own full document.  The complete native source of the article as distributed in the arXiv e-print for this exact version is bundled unchanged as \texttt{sources/208181/source0.tex}.
\begin{lemma}[Polynomial BFC bound]\label{lem:bfc}
Every conjugacy class of a finite group $G$ has size at most
\[
B_N:=(2N+1)^2.
\]
\end{lemma}

\begin{corollary}[Small derived subgroup]\label{cor:derived}
There is an absolute constant $C$ such that every finite group with $\omega(G)=N$ satisfies
\[
\log_2|G'|\le C\bigl(\log_2(N+2)\bigr)^2.
\]
\end{corollary}

\begin{lemma}[Scalar clique credit]\label{lem:scalar-credit}
For every alternating form of rank $\rho$ over $\F_p$, there is a pairwise nonorthogonal set $C$ satisfying
\[
|C|-1\ge \kappa_p\rho-c_p.
\]
\end{lemma}

\begin{equation}\label{eq:central-series}
1=K_0<K_1<\cdots<K_L=P'
\end{equation}

\begin{lemma}[Branch cover bound]\label{lem:branch}
The recursive cover satisfies
\[
\log_2 a(P)\le \max_B\left(\frac{\log_2 p}{2}\sum_{j=0}^{L-1}\rho_j(B)+L\right),
\]
where the maximum is over root-to-leaf branches $B$ and $\rho_j(B)$ denotes the rank encountered at stage $j$ on that branch.
\end{lemma}

\begin{corollary}[Selected-stage composition]\label{cor:selected}
For every subset $E\subseteq\{0,\ldots,L-1\}$,
\[
n\ge\kappa_p\sum_{k\in E}\rho_k-2\kappa_p\!\sum_{\substack{j<k\\ j,k\in E}}t_{j,k}-O\!\left(\kappa_p|E|^2\ell+c_p|E|\right).
\]
\end{corollary}

\begin{corollary}[Expensive stages have small interaction]\label{cor:small-interaction}
There is an absolute constant $C_0$ such that, whenever
\[
\kappa_p\rho_k^2\ge C_0n\ell,
\]
one has, for every $j<k$,
\[
t_{j,k}=O\!\left(\frac{n}{\kappa_p\rho_k}\right),
\]
with an absolute implied constant.
\end{corollary}

\begin{theorem}[Uniform $p$-group bound]\label{thm:pgroup}
Let $P$ range over finite nonabelian $p$-groups and put $n=\omega(P)$. Uniformly in the prime $p$,
\[
\log_2a(P)\le\alpha_p n+O\!\left(\frac{n}{\log(n+2)}\right),
\]
where
\[
\alpha_2=\frac12,\qquad \alpha_3=\frac{\log_2 3}{4},\qquad \alpha_p=\frac{\log_2 p}{p}\quad(p\ge5).
\]
In particular,
\[
\log_2a(P)\le\frac n2+O\!\left(\frac{n}{\log(n+2)}\right),
\]
with strict asymptotic slack for every odd prime.
\end{theorem}

\begin{proof}
Fix an arbitrary root-to-leaf branch of the recursive cover, and use the associated quantities $\rho_k$ and $t_{j,k}$. All estimates below are uniform in the chosen branch, so Lemma~\ref{lem:branch} may be applied at the end. By Lemma~\ref{lem:bfc}, $P$ is a $B$-BFC group with $B=O(n^2)$. By Corollary~\ref{cor:derived},
\[
\log_2|P'|=O((\log n)^2),
\]
so the series \eqref{eq:central-series} has length
\[
L=O\!\left(\frac{(\log n)^2}{\log p}\right).
\]

First suppose $p\le(\log n)^3$. Put $Q=(\log n)^{10}$ and call a stage expensive if $\kappa_p\rho_k\ge n/Q$. The total weighted rank of cheap stages is at most
\[
L\frac nQ=O\!\left(\frac{n}{(\log n)^8}\right).
\]
For an expensive stage, $\rho_k\ge n/(\kappa_pQ)$. Since $p\le(\log n)^3$, we have $\kappa_p=O((\log n)^3)$ and $\ell=O(\log n)$. Consequently
\[
\frac{\kappa_p\rho_k^2}{n\ell}\ge\frac{n}{\kappa_pQ^2\ell}\longrightarrow\infty,
\]
because $\kappa_pQ^2\ell=O((\log n)^{24})$. Thus, for all sufficiently large $n$, the hypothesis of Corollary~\ref{cor:small-interaction} holds uniformly in $p$ and in the branch. Thus $t_{j,k}=O(Q)$ for every interaction entering an expensive stage. There are $O(L^2)$ pairs, so
\[
\kappa_p\sum_{\substack{j<k\\ k\text{ expensive}}}t_{j,k}=o(n).
\]
The anchor loss $O(\kappa_pL^2\ell+c_pL)$ is also $o(n)$. Applying Corollary~\ref{cor:selected} to the expensive stages gives
\[
\sum_{k\text{ expensive}}\kappa_p\rho_k
\le n+O\!\left(\frac{n}{(\log n)^8}+(\log n)^{17}\right).
\]
Adding the cheap stages gives the same bound for $\sum_k\kappa_p\rho_k$. By Lemma~\ref{lem:branch},
\[
\log_2a(P)\le\frac{\log_2p}{2\kappa_p}
\left(n+O\!\left(\frac{n}{(\log n)^8}+(\log n)^{17}\right)\right)+O(L)
=\alpha_p n+O\!\left(\frac{n}{\log(n+2)}\right).
\]

Now suppose $p>(\log n)^3$. Since $P$ is nonabelian, some stage has positive rank; Lemma~\ref{lem:scalar-credit} then gives a clique of size at least $p+1$, so $p<n$. Every positive-rank stage has an individual clique of size at least $(p/2)\rho_k+1$, hence $\rho_k\le2(n-1)/p$. Therefore
\[
\log_2a(P)\le O\!\left(L\frac{n\log p}{p}\right)
=O\!\left(\frac{n(\log n)^2}{p}\right)
=O\!\left(\frac{n}{\log n}\right).
\]
All implied constants are absolute.
\end{proof}

\end{document}
